\DocumentMetadata{
  pdfversion=2.0,
  pdfstandard=ua-2,
  lang=en-US,
  tagging=on,
  tagging-setup={math/setup=mathml-SE,tabsorder=structure}
}
\documentclass[11pt,letterpaper]{article}
\usepackage[margin=0.68in,includefoot]{geometry}
\usepackage{newcomputermodern}
\usepackage{amsmath,amscd,mathtools}
\usepackage{tikz,adjustbox}
\providecommand{\square}{\mdlgwhtsquare}
\usepackage[hidelinks]{hyperref}
\hypersetup{
  pdftitle={Numerical Analysis PhD Exam, May 2014},
  pdfauthor={University of Florida Department of Mathematics},
  pdfsubject={Graduate mathematics examination},
  pdfdisplaydoctitle=true,
  bookmarksopen=true
}
\setcounter{secnumdepth}{0}
\setlength{\parindent}{0pt}
\setlength{\parskip}{0.28em}
\setlength{\emergencystretch}{3em}
\setlength{\leftmargini}{2.15em}
\setlength{\leftmarginii}{2.65em}
\setlength{\leftmarginiii}{3em}
\pagestyle{empty}
\raggedbottom
\allowdisplaybreaks
\NewTaggingSocketPlug{block/list/label}{examproblem}{%
  \tagstructbegin{tag=Lbl}\tagstructend%
  \tagstructbegin{tag=\LBody}%
  \tagstructbegin{tag=\ProblemHeadingTag,title-o={\ProblemHeadingText}}%
  \tagmcbegin{tag=\ProblemHeadingTag,actualtext={\ProblemHeadingText}}%
  #2%
  \tagmcend%
  \tagstructend%
}
\newenvironment{examproblems}[2]{%
  \def\ProblemHeadingTag{#2}%
  \AssignTaggingSocketPlug{block/list/label}{examproblem}%
  \begin{enumerate}%
  \setcounter{enumi}{#1}%
  \renewcommand{\labelenumi}{\textbf{\theenumi.}}%
  \setlength{\topsep}{0.35em}%
  \setlength{\partopsep}{0pt}%
  \setlength{\itemsep}{0.48em}%
  \setlength{\parsep}{0.18em}%
}{\end{enumerate}}
\newenvironment{examsubparts}{%
  \AssignTaggingSocketPlug{block/list/label}{default}%
  \begin{enumerate}%
  \setlength{\topsep}{0.18em}%
  \setlength{\partopsep}{0pt}%
  \setlength{\itemsep}{0.16em}%
  \setlength{\parsep}{0.1em}%
}{\end{enumerate}}
\newenvironment{examinstructions}{%
  \par\begingroup\itshape
}{%
  \par\endgroup\vspace{0.18em}
}
\makeatletter
\renewcommand\subsection{\@startsection{subsection}{2}{0pt}%
  {-1.25ex plus -.25ex minus -.1ex}%
  {0.55ex plus .1ex}%
  {\normalfont\large\bfseries}}
\renewcommand{\@maketitle}{%
  \newpage\null\vskip -1em%
  {\footnotesize\bfseries
    \href{https://www.ufl.edu/}{University of Florida}, \href{https://math.ufl.edu/}{Department of Mathematics}\par}%
  \vskip -0.5em%
  \tagstructbegin{tag=H1,title={Numerical Analysis PhD Exam, May 2014}}%
  \tagpdfparaOff%
  \tagmcbegin{tag=H1}%
    {\LARGE\bfseries\href{https://gma.math.ufl.edu/exams/numerical-analysis-phd/}{Numerical Analysis PhD Exam}, May 2014\par}%
  \tagmcend%
  \tagpdfparaOn%
  \tagstructend%
  \vskip 0.25em%
  \tagmcbegin{artifact}%
    \hrule height 1pt%
  \tagmcend%
  \vskip 1em%
}
\makeatother
\title{Numerical Analysis PhD Exam, May 2014}
\author{}
\date{}
\begin{document}
\maketitle
\thispagestyle{empty}
\begin{examinstructions}
You have 2 hours to answer the following questions. You must show all your work as neatly and clearly as possible and indicate the final answer clearly. You may not use a calculator.
\end{examinstructions}
\begin{examproblems}{0}{H2}
\def\ProblemHeadingText{Problem 1.}
\item
Show that the fixed point equation \(x=f(x)\) with a fixed point~\(\alpha\) can be solved iteratively, if \(f'(\alpha)\ne 1\), by one of the following fixed point iteration formulas:
\begin{examsubparts}
\item[\textnormal{(a)}]
\(x_{n+1}=f(x_n)\)
\item[\textnormal{(b)}]
\(x_{n+1}=f^{-1}(x_n)\). Hint: The following formula for the derivative of an inverse is valid
\[
[f^{-1}]'(x)=\frac{1}{f'(f^{-1}(x))}.
\]
\item[\textnormal{(c)}]
\(x_{n+1}=(x_n+f(x_n))/2\)
\end{examsubparts}
\def\ProblemHeadingText{Problem 2.}
\item
Let \(f(x)\in C^1[a,b]\). Let \(p(x)\) be a polynomial for which
\[
\lVert f'-p\rVert_\infty\le \epsilon
\]
 and define
\[
q(x)=f(a)+\int_a^x p(t)\,dt,\qquad a\le x\le b.
\]
 Show that \(q(x)\) is a polynomial that satisfies
\[
\lVert f-q\rVert_\infty\le \epsilon(b-a).
\]
\def\ProblemHeadingText{Problem 3.}
\item
Consider the integral
\[
\int_0^\pi x^2\cos x\,dx. \tag{1}
\]
\begin{examsubparts}
\item[\textnormal{(a)}]
Consider a quadrature rule of the form
\[
\int_0^\pi x^\alpha f(x)\,dx\approx Af(0)+B\int_0^\pi f(x)\,dx
\]
 where \(\alpha>-1\), \(\alpha\ne 0\) is a parameter. Determine the constants \(A,B\) so that the quadrature formula has degree of exactness one.
\item[\textnormal{(b)}]
Use the formula in part (a) to approximate the integral (1).
\end{examsubparts}
\def\ProblemHeadingText{Problem 4.}
\item
For the function \(f(x)=\ln(1+x)\) for \(x\in[0,1]\), find the minimax approximation polynomial of degree one. Give the exact value of the minimax error.
\def\ProblemHeadingText{Problem 5.}
\item
Assume that you are solving the initial value problem
\[
\begin{aligned}
y'&=f(t,y), && a\le t\le b,\\
y(a)&=\alpha.
\end{aligned}
\]
\begin{examsubparts}
\item[\textnormal{(a)}]
Derive the formula for the global error of the numerical solutions for the ODE problem above obtained via Euler’s method. Hint: The formula is \(M=\lVert Y''\rVert_\infty\):
\[
\lvert Y(t_i)-w_i\rvert<\frac{hM}{2L}\left[e^{L(b-a)}-1\right].
\]
\item[\textnormal{(b)}]
Compute the value of \(M=\lVert Y''\rVert_\infty\) necessary to apply the global error formula above to the specific ODE problem
\[
\begin{aligned}
y'&=\sin(t+2y)+e^t, && 0\le t\le 1,\\
y(0)&=0.
\end{aligned}
\]
\end{examsubparts}
\end{examproblems}
\end{document}
